% !TeX program = lualatex
% Compile twice: lualatex 171.tex
% All references and prose are embedded; no BibTeX database or external inputs.
\documentclass[11pt,a4paper]{article}
\usepackage[margin=24mm,headheight=14pt]{geometry}
\usepackage{fontspec}
\setmainfont{Latin Modern Roman}
\setsansfont{Latin Modern Sans}
\setmonofont{Latin Modern Mono}
\usepackage{amsmath,amssymb,mathtools}
\usepackage{unicode-math}
\setmathfont{Latin Modern Math}
\usepackage{microtype}
\usepackage{enumitem,xurl,needspace}
\usepackage[dvipsnames]{xcolor}
\usepackage{hyperref}
\hypersetup{unicode=true,colorlinks=true,linkcolor=MidnightBlue,urlcolor=MidnightBlue,
 pdftitle={500 Mathematical Problem Statements},pdfauthor={Source-annotated compilation},bookmarksnumbered=true}
\usepackage{bookmark}
\usepackage{tocloft}
\setlength{\cftsecnumwidth}{3em}
\cftsetpnumwidth{2.5em}
\cftsetrmarg{4em}
\usepackage{titlesec}
\titleformat{\section}{\Large\bfseries\raggedright}{\thesection}{0.7em}{}
\usepackage{fancyhdr}
\pagestyle{fancy}\fancyhf{}
\fancyhead[L]{\small\sffamily Mathematical problem statements}
\fancyhead[R]{\small Source edition 22}
\fancyfoot[C]{\thepage}
\setlength{\parindent}{0pt}
\setlength{\parskip}{5pt plus 1pt}
\setlength{\emergencystretch}{3em}
\setcounter{tocdepth}{1}
\setcounter{secnumdepth}{1}
\setlist[itemize]{leftmargin=3em,itemsep=5pt,topsep=4pt,parsep=0pt}
\allowdisplaybreaks
\newcommand{\N}{\mathbb{N}}
\newcommand{\Z}{\mathbb{Z}}
\newcommand{\Q}{\mathbb{Q}}
\newcommand{\R}{\mathbb{R}}
\newcommand{\C}{\mathbb{C}}
\newcommand{\F}{\mathbb{F}}
\newcommand{\A}{\mathbb{A}}
\newcommand{\PP}{\mathbb P}
\newcommand{\E}{\mathbb{E}}
\newcommand{\eps}{\varepsilon}
\newcommand{\dd}{\,\mathrm d}
\newcommand{\sref}[2]{\hyperlink{source:#1:#2}{[S#2]}}
\newcommand{\eref}[2]{\hyperlink{extra:#1:#2}{[E#2]}}
% Turn the next switch off for a shorter reading copy. The quotations preserve
% every exactTarget field, including qualifications omitted from a short summary.
\newif\ifshowcatalogue
\showcataloguetrue
\newcommand{\cataloguescope}[1]{\ifshowcatalogue
 \paragraph{Exact scope in the supplied catalogue.}
 \begin{quote}\small #1\end{quote}\fi}

% Standalone notebook wrapper; original problem section retained verbatim.
\numberwithin{equation}{section}
\hypersetup{pdftitle={Research attempt -- problem 171},
 pdfauthor={Research notebook; source compilation retained}}
\fancyhead[L]{\small\sffamily Research notebook}
\fancyhead[R]{\small\sffamily Problem 171 | 27 September 2026}
\begin{document}
\setcounter{section}{170}
% ================================================================
% problemId: problem.simplicity-of-nontrivial-zeros-of-the-riemann-zeta-function
% source releaseRank: 171; source releaseStatus: open_with_solved_subcases
% exactTarget SHA-256: c7ab74074702f2377cb70d224657ae082e2d221dd47eb7ed0e11076994f2712b
\clearpage
\section{\texorpdfstring{Simplicity of the Nontrivial Zeros of the Riemann Zeta Function}{Simplicity of the Nontrivial Zeros of the Riemann Zeta Function}}\label{171}
\subsection{Definitions and mathematical statement}
Let $\zeta(s)$ be the meromorphic continuation of $\sum_{n\ge1}n^{-s}$ from $\Re s>1$. A nontrivial zero is a zero $\rho$ in $0<\Re\rho<1$. Its multiplicity is the smallest positive integer $m$ such that $\zeta^{(m)}(\rho)\ne0$. The assertion is
\[
 \zeta(\rho)=0,\quad0<\Re\rho<1\quad\Longrightarrow\quad\zeta'(\rho)\ne0.
\]
It concerns every zero in the strip, not only a positive proportion, and does not assume or assert that its real part is $1/2$. Consequently simplicity and the Riemann Hypothesis are distinct targets even though both concern the same zeros.
\par\smallskip\textit{Formulation and concept references:} \sref{171}{1}.
\cataloguescope{Prove that zeta'(rho) is nonzero for every nontrivial zero rho of the Riemann zeta function; equivalently, every nontrivial zero of zeta(s) has multiplicity exactly one.}
\subsection{Short English statement}
Does the zeta function cross zero with multiplicity one at every nontrivial zero, wherever that zero lies in the critical strip?
% BEGIN RESEARCH ADDITION ATTEMPT-171
\subsection{Research attempt}

\paragraph{Current frontier.}
Theorem 1 of Bui--Heath-Brown proves a proportion of at least $19/27$ simple
zeros \emph{assuming RH}; the hypothesis cannot be suppressed
\sref{171}{1}. A 2026 manuscript of Alp\"oge--Furman claims a proportion
$0.6725$ of zeros simple and on the critical line without RH
\eref{171}{1}. That is recorded as a current manuscript claim, not an
independently checked theorem here. Even a verified positive proportion,
or density one, would leave the all-zero assertion unresolved. The attempts
below do not assume RH and do not restrict the zero to the critical line.

\paragraph{Attack route.}
A mollifier route needs absolute control of the exceptional multiplicity,
not merely its average density. A local analytic route can certify one
simple zero at a time, but would need certificates covering the whole
unbounded strip. A third route asks whether integrability of the reciprocal
could rule out repeated zeros without estimating $\zeta'$ individually.
We derive exact criteria for the last two routes and test how much a
mean-value estimate would really have to accomplish.

\paragraph{Attempt: a reciprocal-integrability equivalence.}
Let $\rho$ be any zero strictly inside the critical strip. On a disk
containing no other zero, analytic factorization gives
\[
 \zeta(s)=(s-\rho)^m g(s),\qquad g(\rho)\ne0,
\]
where $m\ge1$. On a smaller closed disk, $|g|$ has positive lower and finite
upper bounds. Consequently, for any real $p>0$,
\begin{equation}\label{eq171-integrable}
 \int_{|s-\rho|<r}|\zeta(s)|^{-p}\,dA(s)<\infty
             \quad\Longleftrightarrow\quad pm<2.
\end{equation}
This follows by polar coordinates: the comparison integral is a constant
multiple of $\int_0^r t^{1-pm}\,dt$. In particular, for $p=1$ the integral
is finite exactly when the zero is simple. At a double zero its divergence
is logarithmic, so the borderline has not been lost.

Thus the target is equivalent to
\begin{equation}\label{eq171-L1}
       1/\zeta\in L^1_{\mathrm{loc}}\bigl(\{0<\Re s<1\},dA\bigr).
\end{equation}
To justify the global wording, a compact subset of the strip meets only
finitely many zeros; away from their small disks the reciprocal is bounded.
The pole of $\zeta$ at $1$ is outside the open region and is not an issue.
Equation \eqref{eq171-L1} is a precise alternative target, not a proved
regularity estimate. Meromorphicity alone allows higher-order poles of the
reciprocal and supplies no such $L^1$ conclusion.

An apparent shortcut is to prove \eqref{eq171-integrable} for every
$p<1$ and let $p\uparrow1$. This also fails: a double zero has finite
local $L^p$ reciprocal norm for every $p<1$. A uniform bound as $p\uparrow1$,
or a direct endpoint estimate, would be required. Conversely, an $L^2$
reciprocal estimate would be too strong: it fails even near a simple zero.
These endpoint tests make the required normalization unambiguous.

\paragraph{Attempt: a finite analytic certificate.}
Let the closed disk $\overline{D(z_0,r)}$ lie in the strip, and set
\[
 a=\zeta(z_0),\qquad b=\zeta'(z_0),\qquad
 M_2\ge\sup_{|s-z_0|\le r}|\zeta''(s)|.
\]
The following strict inequality is a sufficient certificate for exactly one
zero in the disk, counted with multiplicity:
\begin{equation}\label{eq171-rouche}
                    |a|+\tfrac12M_2r^2<|b|r.
\end{equation}
Indeed Taylor's formula along the complex line segment, with integral
remainder, gives
$|\zeta(s)-b(s-z_0)|\le |a|+M_2r^2/2$ on the boundary. Rouch\'e's
 theorem compares it with $b(s-z_0)$, which has exactly one zero there.
The strict inequality implies $b\ne0$, and multiplicities in Rouch\'e's
 theorem show that the enclosed zero is simple. No assertion that the
center itself is a zero was needed.

When $M_2>0$, positive radii can satisfy the scalar inequality only when
$|b|^2>2M_2|a|$; even then the disk must stay in the analytic region and
the chosen $M_2$ must be valid on that same disk. Treating a bound at the
center as a bound on the disk would be circular. An exact toy certificate
is $F(z)=z+z^2$, $z_0=0$, $r=1/4$, $M_2=2$, for which
$M_2r^2/2=1/16<1/4$. The accompanying script checks this rational
inequality. It makes no claim to have certified zeta zeros numerically.

A conceivable proof would combine such disks with a certified count of all
zeros in a finite rectangle and then a theorem providing comparable disks
at every larger height. The first task is finite validation, not the second.
In particular, increasingly many successful computations cannot replace a
uniform lower control of the derivative relative to the remainder.
At a multiple zero, the required certificate correctly becomes impossible
for a disk enclosing just that zero.

\paragraph{Attempt: measuring the multiplicity defect.}
For a finite multiset $\mathcal Z_T$ of nontrivial zeros, counted with their
multiplicities, and any auxiliary function $B$ defined there, put
\[
 \mathcal D_B(T)=\sum_{\rho\in\mathcal Z_T}
                          |1-B(\rho)\zeta'(\rho)|^2.
\]
Every occurrence of a multiple zero contributes exactly one. Therefore
\begin{equation}\label{eq171-defect}
 \mathcal D_B(T)\ge
 \sum_{\substack{\rho\text{ distinct in }\mathcal Z_T\\m(\rho)\ge2}}m(\rho).
\end{equation}
A strict bound $\mathcal D_B(T)<2$ would exclude multiple zeros in that
multiset. A bound $o(\#\mathcal Z_T)$ would not: finitely many, or a sparse
infinite sequence of, multiple zeros could survive. This elementary
integer-valued obstruction explains why a proportion theorem is a
qualitatively different target, independently of RH.

Trying to choose $B(\rho)=1/\zeta'(\rho)$ simply assumes simplicity at the
very points under investigation. A usable choice must be constructed
without that assumption and estimated uniformly. The sum is over every
zero in the strip; replacing it by zeros on the critical line also changes
the target. No auxiliary function with the necessary absolute defect bound
has been found here.

\paragraph{Stress test.}
The equivalence \eqref{eq171-L1} is local in two real dimensions. A
one-dimensional line integral has a different integrability threshold and
cannot be substituted. The strict Rouch\'e inequality excludes boundary
zeros; a nonstrict numerical comparison does not. Neither argument uses
an exchange of a zero sum with an unproved convergent reciprocal moment.

The functional equation does not prevent multiplicity: its reflection
symmetries preserve the order of a zero. Likewise, an entire function with
zeta-like reality symmetries may have repeated zeros. No genericity argument
about arbitrary analytic functions controls this one fixed function.
The current proportion manuscript \eref{171}{1}, even if its advertised
verification succeeds, does not furnish the missing local endpoint or
absolute defect estimate.

\paragraph{Outcome.}
\textbf{Outcome: Exploratory approach with unresolved gap.}

The attempt proves an endpoint reciprocal-integrability equivalence, a
quantitative disk certificate, and an absolute multiplicity-defect test.
These are local analytic tools and diagnostics, with no claim of novelty
for their standard analytic ingredients. What remains is a zeta-specific
estimate making one of the tests hold uniformly throughout the strip.
No zero has been proved multiple, and the simplicity assertion for every
zero has not been established.
% END RESEARCH ADDITION ATTEMPT-171
\subsection{Sources}
\begin{itemize}\raggedright
\item[\textnormal{[S1]}]\hypertarget{source:171:1}{} H. M. Bui and D. R. Heath-Brown, "On simple zeros of the Riemann zeta-function," Bulletin of the London Mathematical Society 45(5) (2013), 953-961.
\url{https://arxiv.org/html/1302.5018}.
{\small\textit{Catalogue formulation source.}}
% BEGIN RESEARCH ADDITION SOURCES-171
\item[\textnormal{[E1]}]\hypertarget{extra:171:1}{} Levent Alp\"oge and
Ralph Furman, \emph{More than two thirds of the zeta zeros are simple and on
the critical line}, arXiv:2608.13637, 13 August 2026, abstract.
\url{https://arxiv.org/abs/2608.13637}.
{\small\textit{Consulted 27 September 2026. Current manuscript claim;
its proof and advertised formal verification have not been audited here.}}
% END RESEARCH ADDITION SOURCES-171
\end{itemize}

\end{document}
